\section{A finite direction alphabet and its exact scope}
\label{sec:finite-direction-alphabet}

The grouped Cauchy transfer permits a different signed residue vector for
 each actual target column.  Consider query rows indexed by an address
 $i\in[G]$ and one of two bit types.  With a fixed invertible matrix $C$, a
 label $(s_g,k_g)\in\{-1,1\}\times\{-1,0,1\}$ at residue $g$ gives scores
\[
 y=Cc,\qquad y'=C\operatorname{diag}(s_g(1+\delta k_g))c.
\]
The two displayed vectors are query row types within one target column;
 they are not two target columns forced to share amplitudes.  Nevertheless,
 arbitrary labels need not satisfy the balanced color inventory required
 by the grouped transfer.

\paragraph{A complete finite lemma.}
Suppose $C\in\mathbb R^{G\times G}$ is invertible and every entry of $C$
 is nonzero.  If $\sigma,\tau\in\{-1,1\}^G$ satisfy
\[
 \min\{d_H(\sigma,\tau),d_H(\sigma,-\tau)\}\leq 1,
\]
then for every sufficiently small $\delta>0$ there are labels from the six
 directions above and a residue vector $c$ giving the strict signs
 $\sigma$ and $\tau$.  A single sufficiently small $\delta$ works for all
 pairs satisfying this condition.

\paragraph{Proof.}
Choose $s\in\{-1,1\}$ so that $\sigma$ and $s\tau$ differ in at most one
 coordinate.  If they agree, take $c=C^{-1}\sigma$ and all labels $(s,0)$.
Otherwise let $j$ be the exceptional coordinate.  Put $(y_0)_j=0$ and
 $(y_0)_i=\sigma_i$ for $i\ne j$, and set $c_0=C^{-1}y_0$.  The only
 potentially exceptional case $G=1$ was already covered by global sign
 reversal, so $y_0\ne0$ and $c_0\ne0$.  Choose $g$ with $(c_0)_g\ne0$.
Since $C_{jg}\ne0$, choose $k=\pm e_g$ so that
\[
 z_j<0\quad\hbox{if }\sigma_j=1,\qquad
 z_j>0\quad\hbox{if }\sigma_j=-1,
 \qquad z=C\operatorname{diag}(k)c_0.
\]
Let $T=C\operatorname{diag}(k)C^{-1}$ and define
\[
 y_\delta=y_0-\frac{\delta z_j}{2}e_j,
 \qquad c_\delta=C^{-1}y_\delta.
\]
Then $\operatorname{sign}(y_\delta)=\sigma$ for every $\delta>0$, while
\[
 C(I+\delta\operatorname{diag}(k))c_\delta
 =y_0+\delta z-\frac{\delta z_j}{2}e_j
  -\frac{\delta^2z_j}{2}Te_j.
\]
At $j$ the leading term is $\delta z_j/2$, which has sign $-\sigma_j$.
At all other coordinates the nonzero constant term has sign $\sigma_i$.
For sufficiently small $\delta$ all desired signs therefore hold, and
 multiplication by $s$ gives $\tau$.  There are finitely many eligible
 pairs, so the minimum of their positive admissible thresholds gives one
 common threshold.  This proves the lemma.

In particular, for $G\leq3$ the Hamming condition holds for every pair.
Consequently the three-address success below needs no special Cauchy
 property and is not evidence for a growing superconstant saving.

\paragraph{Exact rational example.}
The accompanying script and JSON use
\[
 C_{ig}=\left(1-\varepsilon(i+g+2)\right)^{-2},
 \qquad 0\leq i,g<3,\quad
 \varepsilon=10^{-4},\quad\delta=\varepsilon^3.
\]
The JSON contains an explicit rational residue vector, labels, and all
 six nonzero scores for each of the $64$ strict sign patterns.  The checker
 recomputes and checks every stored witness against the original matrix.
This is finite exact computational evidence, with the universality for
 sufficiently small $\delta$ additionally supplied by the preceding proof.
Neither the balanced inventory condition nor a growing-address alphabet
 theorem is established.

\paragraph{A complete chamber enumeration rule.}
For a real matrix $A\in\mathbb R^{r\times k}$ of column rank $k$, every
 strict sign pattern attained by $Ax$ is attained at a solution
\[
 A_Ix=\eta,\qquad |I|=k,\quad\det A_I\ne0,
 \quad\eta\in\{-1,1\}^k.
\]
Indeed, scale a strict witness into the nonempty polyhedron
 $\{x:\sigma_i(Ax)_i\geq1\text{ for every }i\}$.
Starting from any feasible point, if its active rows have rank below $k$,
 choose a nonzero direction annihilating those rows.  Some other row has
 nonzero derivative because $A$ has full column rank.  Orient the direction
 so one inactive constraint decreases, and move to the first new active
 constraint.  Its row is independent of the old active rows because its
 derivative is nonzero.  Repeating gives a point with $k$ independent active
 rows.  This proves the claim even for degenerate arrangements.
It yields complete finite enumeration, unlike a chamber-adjacency search
 that only perturbs extreme rays in independent active directions.

\paragraph{Open scope.}
At four addresses the row pairs $(++),(+-),(-+),(--)$ give the first pair
 outside the nesting explanation.  Exact witnesses obtained for fewer than
 all sign patterns in a degenerate arrangement are only lower bounds on
 coverage.  No impossibility claim follows without a complete feasibility
 certificate.  A compiler saving still requires a family with growing
 address count, growing bit count, all strict patterns, compact exterior
 Cauchy poles, and quantitatively controlled balanced inventories.
